\paragraph{Tasks.} Eight simulator tasks span the regime space: LIBERO-Spatial-1,
LIBERO-Object-1, LIBERO-Goal-1, and LIBERO-Long-1 from LIBERO~\citep{liu2023libero}, plus
Robomimic-Lift-PH, Robomimic-Can-PH, Robomimic-Square-PH, and
Robomimic-Transport-PH~\citep{mandlekar2021}. They range from short, single-object reaches to
long, dual-arm, multi-stage manipulation, with three to four tasks per horizon regime.
LIBERO-Object-1 and Robomimic-Transport-PH are held out for H4, locked before training.

\paragraph{Policies.} We use two architectures to test whether the gap is
architecture-general: Diffusion Policy~\citep{chi2023diffusion} (a 1D conditional U-net
denoiser) and ACT~\citep{zhao2023act} (an action-chunking CVAE transformer). Both use a
ResNet-18 image encoder, proprioception concatenated, action chunk eight, prediction horizon
sixteen, observation horizon two. A run is one (task, seed, architecture) triple: eight tasks
$\times$ three seeds $\times$ two architectures $=$ \NRuns{} runs. Checkpoints are saved at
epochs 10, 25, 50, 80, and 120 (oversampling early and mid training, where the gap is most
informative), for \NCheckpoints{} checkpoints in total. With \NRuns{} runs the protocol has
power above 0.95 to detect a ten-percentage-point gap at the Bonferroni-corrected $\alpha$.

\paragraph{The eight offline metrics.} M1 validation L1 (the baseline); M2 delta-action MSE;
M3 action-distribution entropy; M4 latent Mahalanobis distance to the training distribution;
M5 open-loop replay distance in simulation; M6 inter-seed prediction disagreement; M7
trajectory jerk; M8 action confidence. They are mechanistically distinct on purpose:
prediction error, distributional spread, representation drift, dynamics-aware replay,
ensemble variance, and smoothness all probe different failure modes.

\paragraph{Rollouts.} Twenty closed-loop rollouts per checkpoint on a fixed set of initial
conditions shared across all checkpoints of a task. Deterministic DDIM sampling with a fixed
noise seed removes within-checkpoint stochasticity, so the only variance is across the paired
initial conditions. Success uses each benchmark's published success function.

\paragraph{Analysis.} The analysis code is committed and hashed before training. H1 uses a
paired Wilcoxon on per-run gaps; H2 a mixed-effects likelihood-ratio test; H3 paired Wilcoxon
on per-run Spearman correlations with an inner Bonferroni; H4 a cross-validated ridge with a
leave-one-task-out sensitivity check. A simulation-based power analysis is reported alongside.
